% TOP_PROBLEM: {"problemId": "problem.the-hadamard-matrix-existence-conjecture", "canonicalTitle": "The Hadamard Matrix Existence Conjecture", "releaseRank": 338, "primaryDomain": "cryptography_coding_information_optimization", "primaryDomainLabel": "Cryptography, coding, information & optimization", "displayStatus": "open", "releaseStatus": "open"}
% !TeX program = lualatex
% Definition and sources only; this file is not an LLM solution attempt.
\documentclass[11pt]{article}
\usepackage[margin=25mm]{geometry}
\usepackage{fontspec}
\setmainfont{Latin Modern Roman}
\usepackage{amsmath,amssymb,mathtools}
\usepackage{unicode-math}
\setmathfont{Latin Modern Math}
\usepackage{xurl,hyperref}
\hypersetup{colorlinks=true,urlcolor=blue}
\setcounter{secnumdepth}{0}
\setlength{\parindent}{0pt}
\setlength{\parskip}{5pt}
\setlength{\emergencystretch}{3em}
\begin{document}
\section*{\texorpdfstring{The Hadamard Matrix Existence Conjecture}{The Hadamard Matrix Existence Conjecture}}\label{338}
\subsection{Definitions and mathematical statement}
A Hadamard matrix of order $n$ is a real $n\times n$ matrix $H$ whose entries lie in $\{-1,1\}$ and whose rows are mutually orthogonal. Thus
\[
 HH^{\mathsf T}=nI_n.
\]
The existence conjecture asks, for every positive integer $k$, for such a matrix of order $n=4k$. The scalar product condition is exact; it is not an approximate orthogonality requirement. Interchanging or negating rows and columns preserves the condition but is not part of an additional classification task. The assertion quantifies over every multiple of four, not just infinitely many orders or orders supplied by a particular construction.
\par\smallskip\textit{Formulation and concept references:} \href{https://press.princeton.edu/books/hardcover/9780691119212/hadamard-matrices-and-their-applications}{[S1]}.

\subsection{Short English statement}
Does a square matrix of plus and minus ones with mutually orthogonal rows exist at every order divisible by four?
\subsection{Sources}
\begin{itemize}
\item[S1]K. J. Horadam, Hadamard Matrices and Their Applications, Princeton Univ. Press, 2007.
\url{https://press.princeton.edu/books/hardcover/9780691119212/hadamard-matrices-and-their-applications}.
\textit{Formulation source.}
\end{itemize}
\end{document}
